\section{Research Questions and Methodological Framing}
\label{sec:research_questions}

The seventeen-layer architecture proposed in this paper is not an end in itself but a structured response to a single overarching research question.

\subsection{Main Research Question}
\label{subsec:main_rq}

\begin{quote}
\textbf{How can a computational system autonomously construct a body of knowledge from elementary observables, without inheriting human categories, temporalities, or causal assumptions, such that the resulting knowledge is coherent, verifiable, and capable of offering a genuinely non-anthropocentric perspective on the domain it observes?}
\end{quote}

This question is deliberately formulated to avoid presupposing human epistemic categories. It demands that we specify (i) what constitutes an \emph{elementary observable} (Layer~1), (ii) how structure, function, and temporality can emerge from relations among such observables (Layers~2--4), (iii) how the system can learn to optimize its own organization (Layers~5--7), (iv) how it can generate new ontological categories and reconcile incommensurable frameworks (Layers~8--13), and (v) how it can build self-sustaining worlds and act responsibly at planetary scales (Layers~14--17).

\subsection{Methodological Approach}
\label{subsec:methodology}

The research programme is organized around three interlocking methodological pillars:
\begin{enumerate}
    \item \textbf{Formal specification.} Every layer is defined in the language of category theory, measure theory, dynamical systems, and information theory. This provides a precise, implementation-independent semantics.
    \item \textbf{Falsifiable experimentation.} For each cluster of layers, we formulate hypotheses that can be empirically tested. The reference implementation serves as a testbed for these hypotheses.
    \item \textbf{Iterative refinement.} The \texttt{MetaSpecification} and the evolutionary feedback loops (Layer~6) allow the framework to revise its own definitions of atomicity, coherence, and generativity based on empirical performance.
\end{enumerate}

\subsection{Detailed Research Questions by Cluster}
\label{subsec:detailed_rq}

The \Apeiron{} Framework is not merely a descriptive taxonomy but a research programme designed to be empirically testable and progressively refinable. To anchor this programme, we articulate the following guiding research questions, organized by cluster. The numbering continues across clusters for ease of reference.

\subsubsection{Cluster I: Substrate \& Emergence (Layers 1--4)}
\begin{enumerate}[label=RQ\arabic*:, leftmargin=2.5cm, series=rq]
    \item \textbf{Multi-axial convergence:} Under what conditions do the Boolean, measure-theoretic, categorical, and information-theoretic atomicity frameworks converge on a consistent verdict for a given observable? What are the falsifiable bounds on this convergence?
    \item \textbf{Relational emergence threshold:} What is the minimal relational density (edge-to-node ratio, hyperedge arity) required for the emergence of statistically significant functional motifs?
    \item \textbf{Functional stability:} How does the robustness of a functional entity (measured via counterfactual ablation) scale with the modularity of its underlying relational cluster?
    \item \textbf{Feedback-induced bifurcations:} What are the critical parameter regimes (learning rate $\eta$, noise scale $\sigma$) at which a homeostatic feedback loop transitions into oscillatory or chaotic behavior?
\end{enumerate}

\subsubsection{Cluster II: Adaptive Intelligence (Layers 5--7)}
\begin{enumerate}[label=RQ\arabic*:, leftmargin=2.5cm, resume=rq]
    \item \textbf{Meta-learning transfer:} How does the performance of a meta-learned update rule $\mathcal{H}$ degrade as the distributional distance between training and test task domains increases?
    \item \textbf{Synthesis compression ratio:} What is the relationship between the description length of the global synthesis $\Sigma$ and the predictive accuracy of the lower-layer projections $\Pi_k(\Sigma)$?
    \item \textbf{Guidance stability:} Does the soft guidance operator $\Gamma$ (Layer~7) induce stable attractors in the joint state space of Layers 1--6, and under what conditions does it lead to lock-in?
\end{enumerate}

\subsubsection{Cluster III: Ontological Fluidity (Layers 8--13)}
\begin{enumerate}[label=RQ\arabic*:, leftmargin=2.5cm, resume=rq]
    \item \textbf{Temporal invariant detection:} Can the system reliably identify fixed points of the flux operator $\Delta$ in high-dimensional state spaces, and what are the early-warning signals for impending temporal reconfigurations?
    \item \textbf{Genuine ontological creation:} What is the falsifiable criterion for determining that a newly generated category in Layer~9 is genuinely irreducible to combinations of existing categories?
    \item \textbf{Coherence-revision trade-off:} How does the global coherence threshold $\mathcal{K}_{\text{min}}$ affect the rate of ontological revision and the long-term stability of the architecture?
    \item \textbf{Ontogenesis precursors:} Do persistent homology and critical slowing down provide statistically reliable precursors to ontogenetic events in Layer~13?
\end{enumerate}

\subsubsection{Cluster IV: World-Making \& Governance (Layers 14--17)}
\begin{enumerate}[label=RQ\arabic*:, leftmargin=2.5cm, resume=rq]
    \item \textbf{Plurality preservation:} What is the minimal Gromov-Hausdorff distance $\varepsilon$ required to maintain distinct world-models in a pluriverse without collapse into a single dominant ontology?
    \item \textbf{Reversibility guarantees:} What classes of morphogenetic actions admit verifiable $\varepsilon$-reversibility, and what are the computational complexity bounds for certifying such reversibility?
    \item \textbf{Collective transcendence:} Under what network topologies and communication protocols does the collective cognitive operator $\Gamma$ exhibit capacity strictly greater than the maximum individual capacity?
    \item \textbf{Fixed-point convergence:} Does the meta-operator $\mathcal{M}$ of Layer~17 possess a unique stable fixed point, and what are the convergence rates under iterative refinement?
\end{enumerate}

These questions provide a concrete roadmap for empirical and theoretical investigation. Each layer description in the following sections explicitly references the relevant research questions.